\documentclass[11pt]{article}
\usepackage[margin=0.75in]{geometry}
\usepackage{fontspec}
\usepackage{setspace}
\usepackage{hyperref}
\hypersetup{colorlinks=true, urlcolor=blue, linkcolor=black}
\pagestyle{empty}
\setlength{\parskip}{0.45em}
\setlength{\parindent}{0pt}
\setstretch{0.94}
\setmainfont{Times New Roman}

\begin{document}

\begin{flushleft}
Ke Cheng HSU (Clio) \\
Department of Economics, The Hong Kong University of Science and Technology \\
Clear Water Bay, Kowloon, Hong Kong $\cdot$ +852 6840 9753 \\
\href{mailto:khsu@connect.ust.hk}{khsu@connect.ust.hk} $\cdot$ \href{https://cliohsu1997.github.io/}{cliohsu1997.github.io}
\end{flushleft}

10 November 2026

\vspace{0.15em}
Faculty Search Committee \\
Institute of Economics, Academia Sinica \\
128 Academia Road, Section 2, Nangang District \\
Taipei 115201, Taiwan

\vspace{0.15em}
Dear Members of the Search Committee,

I am writing to apply for the Assistant Research Fellow position at the Institute of Economics, Academia Sinica. I am a Ph.D.\ candidate in Economics at the Hong Kong University of Science and Technology, expecting to graduate in June 2027 and able to join in Summer 2027. My fields are Empirical Industrial Organization and Digital Economics. I grew up in Taiwan and worked as a research assistant at National Taiwan University before beginning my doctoral studies in Hong Kong, and I hope to return home to build my long-term research career at Academia Sinica.

My research uses both theoretical models and structural estimation, which fits the industrial organization group at the Institute. Few junior fellows work on empirical industrial organization in product and platform markets, while the group has strong theorists working on information, platforms, and market design. Working alongside those colleagues is the environment I want, because a theoretical idea on search, ratings, or contract design can be tested in the data, and an empirical finding from an online market can suggest the next model.

My job market paper, ``Designing Return Policy in E-Commerce: A Trade-off between Purchase and Return Efficiency,'' studies whether a return policy encourages efficient returns and purchases. In Taiwan and other online retail markets, easy returns let consumers inspect goods at home, yet frequent returns leave sellers with round-trip shipping and handling costs and lead them to ask whether returns should be made stricter. A stricter policy curbs inefficient returns after delivery, but the lower return rate makes a higher product price more attractive for the seller and can exclude efficient purchases. Using transaction data from a Taobao apparel retailer, I estimate a structural model of purchase, return, and whether to buy return-shipping insurance under monopoly pricing across six categories. Automatic coverage of return shipping performs better when the margin, net scrap value, and uncertainty are high, because it expands efficient purchases. Optional coverage performs better when those attributes are low, because it curbs excessive returns. A menu that pairs a discounted price with a medium restocking fee balances the two margins and performs better than both insurance policies for most categories.

A second working paper, ``Product Design with Heterogeneous Search Costs'' (with Pak Hung Au), studies how firms choose product design when consumers differ in search costs. Search costs shape which products consumers search. Consumers with low search costs seek niche products, and consumers with high search costs seek standardized products. A fall in search costs can raise prices when this shift in which product consumers search outweighs the effect that encourages more search. I also have two ongoing projects. With Kohei Kawaguchi and Kei Kawai, I use Japanese bookstore data to study how resale price maintenance encourages retailers to hold inventory in low-demand states. However, unsold copies cannot be cleared by cutting the price, so they are returned and discarded at an environmental cost. With Ce Bian and Liang Zhao, I build a theoretical model of pass-through when products can be returned, and show that a tax can lower the equilibrium price and raise welfare by reducing the return rate.

I submit my curriculum vitae and job market paper with this application. Pak Hung Au, Kohei Kawaguchi, and Kei Kawai will provide letters of recommendation. You can reach me at khsu@connect.ust.hk or +852 6840 9753.

\vspace{0.35em}
Sincerely,

\vspace{0.35em}
Ke Cheng HSU \\
Ph.D.\ Candidate, Department of Economics, HKUST

\end{document}
